Potholes cause accidents and damage vehicles, yet municipal bodies still find most of them through complaints and manual surveys, and navigation applications choose routes with no information about the road surface. This project rebuilds an earlier segmentation-based prototype as a system that detects potholes in forward-facing road video, places them on a map and compares routes by the potholes expected along them.

Detection is single-class box detection with YOLOv8s. We prepared RDD2022 and HRP4K as one pothole class, grouped near-duplicate images by perceptual hash so that no validation or test image has a training image within 8 bits, and kept 203 images from a third dataset that labels manholes for evaluation only. Thresholds were chosen on validation data, and confidence intervals come from a group bootstrap.

The deployed baseline, trained on RDD2022 at 640 px, reached mAP\textsubscript{50} 0.497 [0.443, 0.555] on the held-out RDD2022 test split. An experimental candidate fine-tuned on RDD2022 and HRP4K at 1280 px raised mAP\textsubscript{50} on the HRP4K test split from 0.018 to 0.454, a gain that comes from training on that data. On RDD2022 its mAP\textsubscript{50} was 0.0445 lower (paired interval [$-$0.0864, $-$0.0059]), and it did not demonstrate fewer false positives or less confusion with manholes, so the baseline remains deployed. Our own recordings have no labels and were compared only qualitatively.

The system covers video and GPS ingestion, a motion-aware tracker, merging of sightings into mapped potholes, snapping to an OpenStreetMap road graph, a Dijkstra route planner whose cost adds expected potholes to distance, a FastAPI backend with SQLite storage and a five-tab browser interface. All 1,101 automated tests pass; the interface has no behaviour tests. For each road segment the system records a pothole count, surveyed length and rate per 100 m, but no severity or condition grade. Route comparison is demonstrated on a synthetic session, because our recordings cover only about 1,020 m of road.
